\section{Setting, terms and rules}\label{sec:setting}

% Figure: where the three modules sit. Every label restates Sections 3-5; no number here is new.
\begin{figure}[!tb]
\centering
\footnotesize
\begin{tikzpicture}[
  box/.style={draw, rounded corners=2pt, align=center, inner sep=4pt, text width=#1, minimum height=0.95cm},
  con/.style={box=#1, fill=blue!8},
  cfm/.style={box=#1, fill=green!12},
  cmp/.style={box=#1, fill=orange!16},
  plain/.style={box=#1, fill=white},
  arr/.style={-{Stealth[length=1.8mm]}, thick},
  darr/.style={{Stealth[length=1.8mm]}-{Stealth[length=1.8mm]}, thick},
  lab/.style={font=\scriptsize, align=center, inner sep=1.5pt, text=black!75}]
% row 1: what starts a job
\node[plain=2.5cm] (jobs)   at (1.2,0)    {job list: task $\times$ arm $\times$ repeat};
\node[con=3.9cm]   (runner) at (5.6,0)    {\textbf{runner} (trusted): starts agent runs only if eight recorded hashes match};
% row 2: the job
\node[lab, anchor=south] (ctrlab) at (5.05,-1.93) {one fresh container per job};
\node[plain=2.25cm] (agent) at (3.6,-2.55) {agent: framework and bundle};
\node[plain=2.25cm] (repo)  at (6.5,-2.55) {task repository and its tests};
\node[lab, anchor=north, text width=5.5cm] (restr) at (5.05,-3.2)
  {no network, no capabilities, seccomp filter, read-only root, unprivileged user; CPU, memory, process and file caps};
\begin{scope}[on background layer]
\node[draw, dashed, rounded corners=3pt, fill=blue!8, inner xsep=7pt, inner ysep=4pt, fit=(ctrlab)(agent)(repo)(restr)] (ctr) {};
\end{scope}
\node[con=2.45cm] (fwd) at (11.55,-2.55) {\textbf{forwarder}: chat requests only, allow-listed fields, approved template};
\node[plain=1.75cm] (srv) at (14.75,-2.55) {model server on one GPU};
% row 3: what reads the outputs
\node[cmp=5.5cm] (cmpn) at (5.05,-6.0)  {\textbf{compare}: candidate against baseline on the same tasks and repeats; repository clusters; frozen pass rule};
\node[cfm=5.2cm] (cfmn) at (13.15,-6.0) {\textbf{conform}: eight checks of requests and rendered prompts against the scorer path, down to token ids};
% arrows
\draw[arr] (jobs) -- (runner);
\draw[arr] (runner.south) -- node[lab, right, xshift=1pt] {starts, grades} (runner.south |- ctr.north);
\draw[darr] (repo.east) -- (fwd.west);
\path (ctr.east |- repo.east) -- node[lab, above] {per-job} node[lab, below] {Unix socket} (fwd.west);
\draw[arr] (ctr.south) -- node[lab, right, xshift=1pt] {verdict rows} (ctr.south |- cmpn.north);
\draw[arr] (fwd.south) -- node[lab, left, xshift=-1pt] {requests} (fwd.south |- cfmn.north);
\draw[arr] (srv.south) -- node[lab, right, xshift=1pt, text width=1.3cm] {rendered prompts} (srv.south |- cfmn.north);
% module tag
\node[lab, anchor=south west, text=blue!50!black] at ([yshift=1pt]runner.north west) {\textbf{contain}};
\end{tikzpicture}
\caption{The kit: where each module sits.}
\label{fig:overview}
\end{figure}


\subsection{The competition and its scorer}

A submission to the competition is an \emph{agent bundle}: a directory of configuration files for the google-adk agent
framework, which names the agents, their prompts, their tools, their sampling settings and their limits. The base
model is not part of the bundle: the scorer loads the bundle into the competition's harness, serves one Gemma 4 model
with vLLM on four L4 GPUs, and runs the agent on hidden tasks one after another. A task is an issue of a public Python
repository at a fixed commit. The agent works in a sandbox with the repository checked out, and ends by submitting a
patch. The task counts as solved if hidden tests pass after the patch is applied. The whole scoring run has a time limit
of 12 hours, so a slow agent loses tasks at the end. A team may submit once per day, and the public leaderboard shows the
solved share on 58 of the hidden tasks.

Two properties of this setting drive the rest of the paper. First, one hosted score is a draw with a standard
deviation near 1.5 tasks of 58 \citep{afshar2026census,mobassir2026toosmall}, and there is one draw per day. Second,
the scorer's configuration is public in a starter notebook \citep{holbrook2026starter}, but a local stack can differ
from it in ways that do not stop a run (\cref{sec:conform}).

\subsection{Models and software}

\paragraph{Stack.} Gemma 4 E4B, 4-bit quantization-aware checkpoint (Kaggle model version 2); vLLM 0.19.1; google-adk
1.36.1; the competition's harness packages. The experiments of \cref{sec:big} use the 31B model, also as a 4-bit
quantization-aware checkpoint. The serving measurements of \cref{sec:serving} add the 12B model. We write 31B, 12B and
E4B for the three models. E4B is one of the two models that are eligible for the small-model prize of the competition,
and it is the model of the kit's main use.

\subsection{Machines}

Agent runs were made on three kinds of machine, and every released row names its kind.

\paragraph{Emulated latency.} The first experiments with the 31B model ran on a rented cloud machine with one large
GPU. Its server answers faster than the scorer's four L4 GPUs, so each model call was padded to the duration that a
latency model predicts for the scorer's hardware: $0.15 + 0.032\,c + 0.0007\,n$ seconds for a call that generates $c$
tokens after $n$ new prompt tokens. We fitted this model to public call logs from real four-L4 hardware
\citep{moshegolderson2026freshswe}. Several jobs ran at the same time, and the load differed between
blocks; this is why arms are compared within a block wherever possible. \Cref{sec:big-real} compares the emulation with
real hardware.

\paragraph{Four L4 GPUs.} Kaggle notebooks with four L4 GPUs are the hardware type of the scorer. Jobs ran there with the
server flags of the starter notebook and without padding. Notebook hours are limited by a weekly quota, so
these blocks are small.

\paragraph{Workstation.} Experiments with E4B from \cref{sec:parity} on ran on a shared single-GPU workstation
(RTX 5090, 32 GB), with eight jobs at a time and each job in a container of the runner of \cref{sec:contain}. This is
the machine for which the kit was built; it finishes about 250 runs per hour.

\subsection{Tasks}

\paragraph{Development tasks.} The competition distributes development tasks with its harness. We used 70 of them,
from four repositories: a first set of 30 and 40 further tasks.
Competition data may not be redistributed, so the released rows name these tasks by opaque codes
(\texttt{dev-001} and so on) and give only their repository.

\paragraph{Berth.} Tasks for the work on E4B come from Berth \citep{patel2026berth}: 579 of its 653 tasks load in the
harness and pass our validation (the reference patch passes, an empty patch fails). The released rows keep the Berth
task identifiers. Four task sets were drawn from these 579 tasks: a screen set S of 48 tasks, a small panel C of 72 tasks in 18 repositories, a large panel of 200 tasks
in 20 repositories that contains C, and the case-study panel of 231 tasks in 15 repositories. We write H for the 128
tasks of the large panel that are outside C. S, the large panel and the case-study panel share no task. The two panels were frozen before any outcome on them
existed (\cref{sec:compare}). S contains the 24 tasks of the parity runs of \cref{sec:parity}, on which results of the
baseline were known before the screens ran; \cref{sec:programmes} comes back to this.

\subsection{Terms}

A \emph{run} is one agent job on one task: the harness prepares the sandbox, the agent works until it submits or is
stopped, and the harness grades the patch. A run is \emph{solved} or not; a run that ends without a patch or by a limit
counts as not solved. An \emph{arm} is one agent configuration, that is, one bundle with its settings. A \emph{block}
is one session on one machine in which one or more arms ran over one task list in a single job queue; a block with
several arms interleaves them, so they share the load and the time of day. A \emph{repeat} is one of several runs of
the same arm on the same task within a block. \emph{Points} are percentage points of the solve rate. A
\emph{task-equivalent} is one task solved in every repeat, or the same number of solved runs spread over tasks.
Arms carry neutral labels (\cref{tab:arms}); block labels name the experiment (\cref{tab:blocks}).

Counts of solved runs come with 95\% Wilson intervals where a single arm is described. Two arms on the same tasks
with one run per task are compared by the exact two-sided sign test on the tasks that only one of them solved. With
repeats, and for every decision about a candidate, the tests of \cref{sec:compare} are used. A null result is worded
``no difference was detected''; none of the experiments below was sized to show that two arms are equal.

\subsection{Threat model}

The adversary is the code the agent writes. It may try to read host files or other jobs, reach the network, exhaust
host resources, forge its verdict, or outlive its job. Kernel exploits are out of scope.

\subsection{Who did the work, and what is withheld}

\paragraph{Who did the work.} AI agents wrote the kit's code, attacks and analyses (builders). Separate AI agents
(verifiers) re-implemented or re-ran the work adversarially. A lead AI session froze the protocols; a human owner
approved isolation, machine time and release. ``Verified'' below means this process, not human review. The same holds
for the experiments that this version adds: AI agents built the arms, ran the blocks and wrote the analysis scripts
and this text, and the owner approved each hosted submission and each release.

\paragraph{What is withheld.} The competition is still running. The E4B agent bundle that serves as the baseline of
\cref{sec:small,sec:case,sec:programmes}, the changed bundles derived from it, their prompts, the fine-tuned adapter
of \cref{sec:small-first} and its training data are withheld for now; the release of the bundle and of the request
dumps of the case study is planned for 2026-12-03. The arms are therefore described by the kind of change and not by
their text (\cref{tab:arms}). What is released is the outcome and the measurements of every run. Two 31B arms are
public bundles run unchanged: \arm{BF5} is the public Budget-Fit V5 configuration, the most-run public
configuration in the hosted-score census of \citet{afshar2026census}, and \arm{ID24} is a bundle of the public collection of \citet{idris2026runs}.
